\documentclass[10pt]{article}
\usepackage[margin=20mm]{geometry}
\usepackage{amsmath,amssymb,fontspec,unicode-math,xurl,hyperref}
\IfFontExistsTF{DejaVu Serif}{\setmainfont{DejaVu Serif}}{\setmainfont{Latin Modern Roman}}
\setmathfont{latinmodern-math.otf}
\setlength{\parindent}{0pt}\setlength{\parskip}{6pt}\sloppy
\begin{document}
\section*{0067 | Scientific verification report}
Position: When Incentives Backfire, Data Stops Being Human\par
Sebastin Santy, Prasanta Bhattacharya, Manoel Horta Ribeiro, Kelsey Allen and Sewoong Oh. ICML 2025, PMLR 267:82151-82165. Reviewed version: retained publisher-linked 15-page PDF. Main argument pp.1-9; references pp.10-15. Page numbers below refer to this PDF unless another source is named.\par
Scientific decision pending; no human approval or author adoption claimed.\par
\subsection*{Outcome}
Supported severity: Minor. One historical experiment is summarized with an inaccurate alignment of outcome timing. The correction affects the evidentiary wording in Section 5; it does not establish that the proposed motivational mechanism is false. No demonstrated Major issue was found within the inspected claims.\par
\subsection*{Scope and contribution}
The paper proposes aligning data-collection environments with intrinsic motivation, retaining compensation, and exploring games as structured but voluntary collection environments. It presents no new theorem, fitted quantitative model, dataset or controlled experiment. The review assessed the main causal narrative, qualifications, illustrative figures, key primary supporting evidence, public versions and material availability. It did not impose a new-experiment requirement on a normative design proposal.\par
- Sections 1-4: data provenance, the quantity-quality framing and context-dependent definitions of quality. Figure 2 was visually checked and is an illustrative diagram with no measured frontier or numerical optimization claim.\par
- Sections 5-6: reward effects, task fragmentation, LLM-assisted work and the proposed feedback loop. Two central primary sources were checked for what they actually measured.\par
- Sections 7-9: resolution of control, product-integrated collection, trust, games and compensation. Read together with explicit feasibility, bias, privacy and manipulation qualifications.\par
- All scientific main-text pages were read. References were used to locate central evidence; the complete bibliography and every historical/platform claim were not independently audited.\par
\subsection*{Source identity}
The original PDF is 2835052 bytes; SHA-256 `5103a0cffa8d1d9a135df6964fec399019be5f8f84803a1d9e5e6171b582c5f9` matches the source-acquisition record. The title and author list agree with PMLR; expanded middle names in the publisher metadata are identity clarifications. arXiv:2502.07732 lists v2 dated 7 June 2025, after v1 dated 11 February 2025. No later version is listed in the checked record.\par
\subsection*{Supporting checks}
The PMLR related-material section exposes no separate experiment repository or supplement. No such material was required to reproduce a new experiment, because none is reported here. OpenReview was behind browser verification; the CACM website returned HTTP 403, so the author-posted copy was used. Human scientific judgment remains pending.\par
No human review is claimed. Unresolved work consists of the broader proposed causal mechanism, any intervention efficacy and sources outside the focused checks; these are not counted as demonstrated errors.\par
\section*{Problem 1 | Outcome timing in the reward example}
\subsection*{Claim and exact location}
Section 5, p.4, paragraph headed Overjustification Effect, summarizes Lepper, Greene and Nisbett (1973). Its concluding sentence states that ’removing rewards can lead to a decline in both engagement and performance.’ The suspect step assigns a shared post-removal interpretation to two measurements made at different stages.\par
Primary-source endpoint | Measurement stage / locator | What the comparison supports\par
Voluntary drawing time | Later classroom observations; printed p.134 / PDF p.6; 7-14 days after the sessions | Subsequent interest under an unrewarded setting\par
Rated drawing quality | Experimental-session drawings; printed pp.133,135 / PDF pp.5,7 | Immediate performance under the experimental reward conditions\par
The quality ratings were 2.18, 2.85 and 2.69 for expected, unexpected and no-award groups, respectively. They were not post-withdrawal quality measurements. Reading the methods and results together therefore does not support attributing a measured later quality decline to reward removal. This is a documentary endpoint/timing check, not a new statistical reanalysis.\par
\subsection*{Alternative reading and exact impact}
The phrase ’These findings suggest’ can be read as proposing an extrapolation. Under that reading, post-removal performance decline is an untested possibility, not a reported result. The review does not claim that such a later effect is impossible or contradicted by other research. The problem is local evidentiary ambiguity in the illustrative example.\par
Minor is supported because the demonstrated defect changes the interpretation of this one endpoint, rather than invalidating a claimed estimated effect or a formal mechanism used elsewhere. The paper’s design proposal does not depend on a measured post-removal quality coefficient, and its Section 5 discussion explicitly permits beneficial, context-sensitive incentives. A stronger severity would require evidence of damage to the actual central recommendation or its principal empirical basis, which was not demonstrated.\par
\subsection*{Verified correction of the attribution}
Replace the closing inference with a distinction between immediate drawing quality and later voluntary engagement. Treat any effect on later performance as an extrapolation requiring separate support. The primary-source comparison verifies this narrow correction only; no intervention, empirical repair or author adoption was performed.\par
\subsection*{Counterevidence retained}
Section 2 explicitly accepts compensation. Section 5 recognizes informational rewards and calibrated incentives as potentially beneficial. Section 6 permits judicious LLM assistance. Sections 7.2 and 8.2 acknowledge feasibility and trust problems; Section 8.3 discusses bias, privacy and manipulation. Thus neither a universal claim that rewards are harmful nor a proven universal advantage of games is an accurate representation of this paper.\par
\section*{Conclusion, references and evidence appendix}
\subsection*{Other central evidence}
Veselovsky et al. (2025) supports prevalent LLM use in the studied crowdwork task and greater output homogeneity; it also reports better keyword-retention scores for LLM-labeled summaries. These findings concern different dimensions of quality. They do not contradict the reviewed paper’s contextual definition of quality and acknowledgement of useful AI assistance. They do not independently establish the entire proposed causal chain from incentive design through motivational crowding-out to LLM use.\par
Inspected item | Outcome | Repair / remaining work\par
Reward-example timing | Minor evidence-scope mismatch | Narrow attribution correction verified; adoption not requested\par
Quantity-quality and middle-control figures | Scope clarification; conceptual proposals | No quantified Pareto frontier or optimum was tested\par
Crowdwork LLM evidence and nuanced design argument | No Issue Found within checked scope | No causal mediation experiment or full literature audit performed\par
Scientific decision | Pending | publication status and human judgment remain open\par
\subsection*{Consulted primary references}
[R1] Santy et al. (2025), PMLR 267:82151-82165. https://proceedings.mlr.press/v267/santy25a.html . Reviewed source: https://raw.githubusercontent.com/mlresearch/v267/main/assets/santy25a/santy25a.pdf . Version check: https://arxiv.org/abs/2502.07732 (v2).\par
[R2] Lepper, Greene and Nisbett (1973), Undermining Children’s Intrinsic Interest with Extrinsic Reward: A Test of the Overjustification Hypothesis. Journal of Personality and Social Psychology 28(1):129-137. Primary PDF hosted by MIT, retained as `sources/0067/agent\_sources/Lepper\_Greene\_Nisbett\_1973.pdf`; source URL and SHA-256 are in acquisition\_manifest.json.\par
[R3] Veselovsky, Horta Ribeiro, Cozzolino, Gordon, Rothschild and West (2025), Prevalence and Prevention of Large Language Model Use in Crowd Work. Communications of the ACM 68(3):42-47; author-posted publisher-article copy: https://researchdmr.com/files/PrevalencePreventionLLM.pdf .\par
\end{document}
